% CV de base (anglais). Gabarit neutre, 1 page.
% L'assistant remplace les champs [[...]] et les entrées d'exemple avec le vrai CV de l'étudiant,
% puis ne modifie PLUS jamais le préambule ni la mise en page.
\documentclass[10pt,letterpaper]{article}
\usepackage[utf8]{inputenc}
\usepackage[T1]{fontenc}
\usepackage{lmodern}
\usepackage[margin=0.6in]{geometry}
\usepackage{enumitem}
\usepackage[hidelinks]{hyperref}
\usepackage{titlesec}
\pagestyle{empty}
\setlength{\parindent}{0pt}
\titleformat{\section}{\large\bfseries\scshape}{}{0pt}{}[\titlerule]
\titlespacing*{\section}{0pt}{8pt}{4pt}
\setlist[itemize]{leftmargin=1.2em,itemsep=1pt,topsep=2pt,parsep=0pt}
\newcommand{\entry}[4]{\textbf{#1} \hfill #2\\ \textit{#3} \hfill \textit{#4}}

\begin{document}
\begin{center}
{\LARGE\bfseries [[Prénom Nom]]}\\[2pt]
[[Ville]] \,|\, [[téléphone]] \,|\, \href{mailto:[[courriel]]}{[[courriel]]} \,|\, \href{[[lien LinkedIn]]}{LinkedIn} \,|\, \href{[[lien GitHub]]}{GitHub}
\end{center}

\section{Education}
\entry{[[Université]]}{[[Ville]]}{[[Programme, ex. B.Eng. in Software Engineering]]}{[[début -- fin prévue]]}

\section{Experience}
\entry{[[Entreprise]]}{[[Ville]]}{[[Titre du poste]]}{[[dates]]}
\begin{itemize}
  \item [[Réalisation concrète avec un résultat chiffré]]
  \item [[Deuxième réalisation]]
\end{itemize}

\section{Projects}
\entry{[[Nom du projet 1]]}{}{[[Technos]]}{[[date]]}
\begin{itemize}
  \item [[Ce que tu as construit et le résultat]]
  \item [[Deuxième puce]]
\end{itemize}
\entry{[[Nom du projet 2]]}{}{[[Technos]]}{[[date]]}
\begin{itemize}
  \item [[Ce que tu as construit et le résultat]]
\end{itemize}

\section{Technical Skills}
\textbf{Languages:} [[ex. Python, Java, C++, TypeScript]]\\
\textbf{Frameworks \& Tools:} [[ex. React, Node.js, Git, Docker]]\\
\textbf{Other:} [[ex. SQL, Linux, REST APIs]]
\end{document}
